\documentclass[a4paper,12pt,twoside]{memoir}

% Castellano
\usepackage[spanish,es-tabla]{babel}
\selectlanguage{spanish}
\usepackage[utf8]{inputenc}
\usepackage[T1]{fontenc}
\usepackage{lmodern} % Scalable font
\usepackage{microtype}
\usepackage{placeins}

\RequirePackage{booktabs}
\RequirePackage[table]{xcolor}
\RequirePackage{xtab}
\RequirePackage{multirow}

% Links
\PassOptionsToPackage{hyphens}{url}\usepackage[colorlinks]{hyperref}
\hypersetup{
	allcolors = {red}
}

% Ecuaciones
\usepackage{amsmath}

% Rutas de fichero / paquete
\newcommand{\ruta}[1]{{\sffamily #1}}

% Párrafos
\nonzeroparskip

% Huérfanas y viudas
\widowpenalty100000
\clubpenalty100000

% Imágenes

% Comando para insertar una imagen en un lugar concreto.
% Los parámetros son:
% 1 --> Ruta absoluta/relativa de la figura
% 2 --> Texto a pie de figura
% 3 --> Tamaño en tanto por uno relativo al ancho de página
\usepackage{graphicx}
\newcommand{\imagen}[3]{
	\begin{figure}[!h]
		\centering
		\includegraphics[width=#3\textwidth]{#1}
		\caption{#2}\label{fig:#1}
	\end{figure}
	\FloatBarrier
}

% Comando para insertar una imagen sin posición.
% Los parámetros son:
% 1 --> Ruta absoluta/relativa de la figura
% 2 --> Texto a pie de figura
% 3 --> Tamaño en tanto por uno relativo al ancho de página
\newcommand{\imagenflotante}[3]{
	\begin{figure}
		\centering
		\includegraphics[width=#3\textwidth]{#1}
		\caption{#2}\label{fig:#1}
	\end{figure}
}

% El comando \figura nos permite insertar figuras comodamente, y utilizando
% siempre el mismo formato. Los parametros son:
% 1 --> Porcentaje del ancho de página que ocupará la figura (de 0 a 1)
% 2 --> Fichero de la imagen
% 3 --> Texto a pie de imagen
% 4 --> Etiqueta (label) para referencias
% 5 --> Opciones que queramos pasarle al \includegraphics
% 6 --> Opciones de posicionamiento a pasarle a \begin{figure}
\newcommand{\figuraConPosicion}[6]{%
  \setlength{\anchoFloat}{#1\textwidth}%
  \addtolength{\anchoFloat}{-4\fboxsep}%
  \setlength{\anchoFigura}{\anchoFloat}%
  \begin{figure}[#6]
    \begin{center}%
      \Ovalbox{%
        \begin{minipage}{\anchoFloat}%
          \begin{center}%
            \includegraphics[width=\anchoFigura,#5]{#2}%
            \caption{#3}%
            \label{#4}%
          \end{center}%
        \end{minipage}
      }%
    \end{center}%
  \end{figure}%
}

%
% Comando para incluir imágenes en formato apaisado (sin marco).
\newcommand{\figuraApaisadaSinMarco}[5]{%
  \begin{figure}%
    \begin{center}%
    \includegraphics[angle=90,height=#1\textheight,#5]{#2}%
    \caption{#3}%
    \label{#4}%
    \end{center}%
  \end{figure}%
}
% Para las tablas
\newcommand{\otoprule}{\midrule [\heavyrulewidth]}
%
% Nuevo comando para tablas pequeñas (menos de una página).
\newcommand{\tablaSmall}[5]{%
 \begin{table}
  \begin{center}
   \rowcolors {2}{gray!35}{}
   \begin{tabular}{#2}
    \toprule
    #4
    \otoprule
    #5
    \bottomrule
   \end{tabular}
   \caption{#1}
   \label{tabla:#3}
  \end{center}
 \end{table}
}

%
% Nuevo comando para tablas pequeñas (menos de una página).
\newcommand{\tablaSmallSinColores}[5]{%
 \begin{table}[H]
  \begin{center}
   \begin{tabular}{#2}
    \toprule
    #4
    \otoprule
    #5
    \bottomrule
   \end{tabular}
   \caption{#1}
   \label{tabla:#3}
  \end{center}
 \end{table}
}

\newcommand{\tablaApaisadaSmall}[5]{%
\begin{landscape}
  \begin{table}
   \begin{center}
    \rowcolors {2}{gray!35}{}
    \begin{tabular}{#2}
     \toprule
     #4
     \otoprule
     #5
     \bottomrule
    \end{tabular}
    \caption{#1}
    \label{tabla:#3}
   \end{center}
  \end{table}
\end{landscape}
}

%
% Nuevo comando para tablas grandes con cabecera y filas alternas coloreadas en gris.
\newcommand{\tabla}[6]{%
  \begin{center}
    \tablefirsthead{
      \toprule
      #5
      \otoprule
    }
    \tablehead{
      \multicolumn{#3}{l}{\small\sl continúa desde la página anterior}\\
      \toprule
      #5
      \otoprule
    }
    \tabletail{
      \hline
      \multicolumn{#3}{r}{\small\sl continúa en la página siguiente}\\
    }
    \tablelasttail{
      \hline
    }
    \bottomcaption{#1}
    \rowcolors {2}{gray!35}{}
    \begin{xtabular}{#2}
      #6
      \bottomrule
    \end{xtabular}
    \label{tabla:#4}
  \end{center}
}

%
% Nuevo comando para tablas grandes con cabecera.
\newcommand{\tablaSinColores}[6]{%
  \begin{center}
    \tablefirsthead{
      \toprule
      #5
      \otoprule
    }
    \tablehead{
      \multicolumn{#3}{l}{\small\sl continúa desde la página anterior}\\
      \toprule
      #5
      \otoprule
    }
    \tabletail{
      \hline
      \multicolumn{#3}{r}{\small\sl continúa en la página siguiente}\\
    }
    \tablelasttail{
      \hline
    }
    \bottomcaption{#1}
    \begin{xtabular}{#2}
      #6
      \bottomrule
    \end{xtabular}
    \label{tabla:#4}
  \end{center}
}

%
% Nuevo comando para tablas grandes sin cabecera.
\newcommand{\tablaSinCabecera}[5]{%
  \begin{center}
    \tablefirsthead{
      \toprule
    }
    \tablehead{
      \multicolumn{#3}{l}{\small\sl continúa desde la página anterior}\\
      \hline
    }
    \tabletail{
      \hline
      \multicolumn{#3}{r}{\small\sl continúa en la página siguiente}\\
    }
    \tablelasttail{
      \hline
    }
    \bottomcaption{#1}
  \begin{xtabular}{#2}
    #5
   \bottomrule
  \end{xtabular}
  \label{tabla:#4}
  \end{center}
}



\definecolor{cgoLight}{HTML}{EEEEEE}
\definecolor{cgoExtralight}{HTML}{FFFFFF}

%
% Nuevo comando para tablas grandes sin cabecera.
\newcommand{\tablaSinCabeceraConBandas}[5]{%
  \begin{center}
    \tablefirsthead{
      \toprule
    }
    \tablehead{
      \multicolumn{#3}{l}{\small\sl continúa desde la página anterior}\\
      \hline
    }
    \tabletail{
      \hline
      \multicolumn{#3}{r}{\small\sl continúa en la página siguiente}\\
    }
    \tablelasttail{
      \hline
    }
    \bottomcaption{#1}
    \rowcolors[]{1}{cgoExtralight}{cgoLight}

  \begin{xtabular}{#2}
    #5
   \bottomrule
  \end{xtabular}
  \label{tabla:#4}
  \end{center}
}



\graphicspath{ {./img/} }

% Capítulos
\chapterstyle{bianchi}
\newcommand{\capitulo}[2]{
	\setcounter{chapter}{#1}
	\setcounter{section}{0}
	\setcounter{figure}{0}
	\setcounter{table}{0}
	\chapter*{\thechapter.\enskip #2}
	\addcontentsline{toc}{chapter}{\thechapter.\enskip #2}
	\markboth{#2}{#2}
}

% Apéndices
\renewcommand{\appendixname}{Apéndice}
\renewcommand*\cftappendixname{\appendixname}

\newcommand{\apendice}[1]{
	%\renewcommand{\thechapter}{A}
	\chapter{#1}
}

\renewcommand*\cftappendixname{\appendixname\ }

% Formato de portada
\makeatletter
\usepackage{xcolor}
\newcommand{\tutor}[1]{\def\@tutor{#1}}
\newcommand{\course}[1]{\def\@course{#1}}
\definecolor{cpardoBox}{HTML}{E6E6FF}
\def\maketitle{
  \null
  \thispagestyle{empty}
  % Cabecera ----------------
\noindent\includegraphics[width=\textwidth]{cabecera}\vspace{1cm}%
  \vfill
  % Título proyecto y escudo informática ----------------
  \colorbox{cpardoBox}{%
    \begin{minipage}{.8\textwidth}
      \vspace{.5cm}\Large
      \begin{center}
      \textbf{TFG del Grado en Ingeniería Informática}\vspace{.6cm}\\
      \textbf{\LARGE\@title{}}
      \end{center}
      \vspace{.2cm}
    \end{minipage}

  }%
  \hfill\begin{minipage}{.20\textwidth}
    \includegraphics[width=\textwidth]{escudoInfor}
  \end{minipage}
  \vfill
  % Datos de alumno, curso y tutores ------------------
  \begin{center}%
  {%
    \noindent\LARGE
    Presentado por \@author{}\\ 
    en Universidad de Burgos --- \@date{}\\
    Tutor: \@tutor{}\\
  }%
  \end{center}%
  \null
  \cleardoublepage
  }
\makeatother

\newcommand{\nombre}{Nombre del alumno} %%% cambio de comando

% Datos de portada
\title{título del TFG}
\author{\nombre}
\tutor{nombre tutor}
\date{\today}

\begin{document}

\maketitle


\newpage\null\thispagestyle{empty}\newpage


%%%%%%%%%%%%%%%%%%%%%%%%%%%%%%%%%%%%%%%%%%%%%%%%%%%%%%%%%%%%%%%%%%%%%%%%%%%%%%%%%%%%%%%%
\thispagestyle{empty}


\noindent\includegraphics[width=\textwidth]{cabecera}\vspace{1cm}

\noindent D. nombre tutor, profesor del departamento de nombre departamento, área de nombre área.

\noindent Expone:

\noindent Que el alumno D. \nombre, con DNI dni, ha realizado el Trabajo final de Grado en Ingeniería Informática titulado título de TFG. 

\noindent Y que dicho trabajo ha sido realizado por el alumno bajo la dirección del que suscribe, en virtud de lo cual se autoriza su presentación y defensa.

\begin{center} %\large
En Burgos, {\large \today}
\end{center}

\vfill\vfill\vfill

% Author and supervisor
\begin{minipage}{0.45\textwidth}
\begin{flushleft} %\large
Vº. Bº. del Tutor:\\[2cm]
D. nombre tutor
\end{flushleft}
\end{minipage}
\hfill
\begin{minipage}{0.45\textwidth}
\begin{flushleft} %\large
Vº. Bº. del co-tutor:\\[2cm]
D. nombre co-tutor
\end{flushleft}
\end{minipage}
\hfill

\vfill

% para casos con solo un tutor comentar lo anterior
% y descomentar lo siguiente
%Vº. Bº. del Tutor:\\[2cm]
%D. nombre tutor


\newpage\null\thispagestyle{empty}\newpage




\frontmatter

% Abstract en castellano
\renewcommand*\abstractname{Resumen}
\begin{abstract}
En este primer apartado se hace una \textbf{breve} presentación del tema que se aborda en el proyecto.
\end{abstract}

\renewcommand*\abstractname{Descriptores}
\begin{abstract}
Palabras separadas por comas que identifiquen el contenido del proyecto Ej: servidor web, buscador de vuelos, android \ldots
\end{abstract}

\clearpage

% Abstract en inglés
\renewcommand*\abstractname{Abstract}
\begin{abstract}
A \textbf{brief} presentation of the topic addressed in the project.
\end{abstract}

\renewcommand*\abstractname{Keywords}
\begin{abstract}
keywords separated by commas.
\end{abstract}

\clearpage

% Indices
\tableofcontents

\clearpage

\listoffigures

\clearpage

\listoftables
\clearpage

\mainmatter
\capitulo{1}{Introducción}

Evaluar las tareas de programación que entregan los estudiantes no solo implica verificar que el código funcione correctamente, sino también valorar aspectos como la calidad del diseño, la claridad, la eficiencia y el cumplimiento de criterios específicos establecidos por el docente.

Tradicionalmente, la corrección manual de estas tareas es una labor que demanda una gran inversión de tiempo y esfuerzo por parte del profesorado, especialmente en cursos con gran número de estudiantes. Esta situación limita la capacidad de proporcionar una retroalimentación rápida, detallada y personalizada, aspectos fundamentales para el aprendizaje eficaz.

En los últimos años, el avance en modelos de lenguaje a gran escala (\textit{LLM}) ha abierto nuevas posibilidades en el campo de la docencia. Estos modelos tienen la capacidad de procesar y comprender tanto texto como código, lo que los convierte en herramientas prometedoras para asistir en la corrección de tareas, ofreciendo comentarios elaborados que van más allá de una simple validación funcional.

Este proyecto se centra en el desarrollo de una aplicación web que integra el potencial de los \textit{LLM} para realizar evaluaciones automatizadas basadas en criterios definidos por el usuario, mediante la incorporación de rúbricas personalizadas y \textit{prompts} específicos diseñados por los usuarios. La aplicación busca facilitar la labor docente, permitiendo una gestión eficiente de las correcciones y mejorando la calidad de la retroalimentación brindada a los estudiantes.

De esta forma, la herramienta propuesta contribuye a modernizar la enseñanza de la programación, promoviendo una evaluación más ágil, objetiva y enriquecedora, que favorezca tanto al profesorado como al alumnado, sin sustituir la necesaria supervisión humana en este proceso.

\section{Estructura de la memoria}

La memoria sigue la siguiente estructura:

\begin{itemize}
    \item \textbf{Introducción:} breve descripción del problema a resolver y la solución propuesta.
    \item \textbf{Objetivos del proyecto:} exposición de los objetivos que persigue el proyecto.
    \item \textbf{Conceptos teóricos:} breve explicación de los conceptos teóricos.
    \item \textbf{Técnicas y herramientas:} listado de técnicas metodológicas y herramientas utilizadas.
    \item \textbf{Aspectos relevantes del desarrollo:} exposición de aspectos destacables que tuvieron lugar durante la realización del proyecto.
    \item \textbf{Trabajos relacionados:} pequeño resumen de los trabajos y proyectos ya realizados en el campo del proyecto en curso.
    \item \textbf{Conclusiones y líneas de trabajo futuras:} conclusiones obtenidas tras la realización del proyecto.
\end{itemize}

Junto a la memoria se proporcionan los siguientes anexos:

\begin{itemize}
    \item \textbf{Plan del proyecto software:} planificación temporal y estudio de viabilidad del proyecto.
    \item \textbf{Especificación de requisitos del software:} se describe la fase de análisis; los objetivos generales, el catálogo de requisitos del sistema y la especificación de requisitos funcionales y no funcionales.
    \item \textbf{Especificación de diseño:} se describe la fase de diseño; el ámbito del software, el diseño de datos, el diseño procedimental y el diseño arquitectónico.
    \item \textbf{Manual del programador:} recoge los aspectos más relevantes relacionados con el código fuente.
    \item \textbf{Manual de usuario:} guía para el correcto manejo de la aplicación IAGScore.
\end{itemize}

Además, se incluyen los siguientes materiales complementarios accesibles a través de Internet:

\begin{itemize}
    \item \textbf{Repositorio:} el código fuente y documentación están disponibles en el repositorio oficial de \href{https://github.com/pac1006/IAGScore}{IAGScore en GitHub}.
    \item \textbf{Gestión del proyecto:} realizada con \href{https://zube.io/miorganizacion/iagscore/w/workspace-1/kanban}{Zube}.
    \item \textbf{Vídeo de presentación:} \href{https://universidaddeburgos-my.sharepoint.com/:v:/g/personal/pac1006_alu_ubu_es/EfB4_46cvc1Ltzs4VFGbCFkBWj3Dm_AiB9mt6hpTRJkKeQ?nav=eyJyZWZlcnJhbEluZm8iOnsicmVmZXJyYWxBcHAiOiJPbmVEcml2ZUZvckJ1c2luZXNzIiwicmVmZXJyYWxBcHBQbGF0Zm9ybSI6IldlYiIsInJlZmVycmFsTW9kZSI6InZpZXciLCJyZWZlcnJhbFZpZXciOiJNeUZpbGVzTGlua0NvcHkifX0&e=3IHW2I}{presentación} general del proyecto y sus objetivos.
    \item \textbf{Vídeo de demostración:} \href{https://universidaddeburgos-my.sharepoint.com/:v:/g/personal/pac1006_alu_ubu_es/ETGJTbAviu5NrIwAoXvITHYBXlwclYbDOcRTdR0pd7QDdw?nav=eyJyZWZlcnJhbEluZm8iOnsicmVmZXJyYWxBcHAiOiJPbmVEcml2ZUZvckJ1c2luZXNzIiwicmVmZXJyYWxBcHBQbGF0Zm9ybSI6IldlYiIsInJlZmVycmFsTW9kZSI6InZpZXciLCJyZWZlcnJhbFZpZXciOiJNeUZpbGVzTGlua0NvcHkifX0&e=vbhU2M}{demostración} funcional de la aplicación en uso.
\end{itemize}

\capitulo{2}{Objetivos del proyecto}

A continuación se detallan los objetivos que motivan la realización de este proyecto:

\section{Objetivos generales}\label{objetivos_generales}

\begin{itemize}
    \item Desarrollar un sistema \textit{software} que permita automatizar la evaluación de tareas de programación en cualquier lenguaje.
    \item Investigar la posibilidad de hacer uso de diferentes modelos de lenguaje avanzado para la evaluación de estas tareas.
    \item Comprobar la posibilidad de delegar totalmente en estos modelos para la realización de estas tareas.
    \item Implementar operaciones e interacción con la aplicación:
    \begin{itemize}
        \item Permitir a los usuarios crear una cuenta.
        \item Permitir al usuario importar documentos en formato \textit{Markdown} con los contenidos de sus propias rúbricas.
        \item Permitir a los usuarios diseñar o definir los \textit{prompts} con los que interactuarán con los modelos.
        \item Permitir a los usuarios configurar diferentes parámetros del modelo que evaluará las tareas.
        \item Permitir a los usuarios ejecutar el modelo para que realice la evaluación.
        \item Permitir a los usuarios consultar el reaultado de la ejecución de la evaluación.
        \item Permitir a los usuarios consultar datos sobre el modelo utilizado.
    \end{itemize}
\end{itemize}

\section{Objetivos técnicos}\label{objetivos_tecnicos}
\begin{itemize}
    \item Crear una aplicación web utilizando python y el \textit{framework} \textbf{Django}
    \item Utilizar la arquitectura MVT(Model - View - Template) de Django.
    \item Utilizar una base de datos \textbf{PostgreSQL} para el almacenamiento de datos.
    \item Utilizar \textbf{Ollama} para la gestión de los modelos de lenguaje.
    \item Utilizar \textbf{Celery} y \textbf{Redis} para la gestión de tareas asíncronas.
    \item Utilizar \textbf{Docker} para la creación de contenedores y facilitar el despliegue del proyecto.
    \item Hacer uso de \textbf{Sphinx} para la generación de la documentación del proyecto.
    \item Gestionar el proyecto utilizando la herramienta de gestión de proyectos \textbf{Zube} para la planificación y seguimiento del proyecto.
    \item Realizar test con una cobertuta que garanticen la calidad del producto final.
    \item Hacer uso del sistema de control de versiones \textbf{GIT} distribuido junto con la plataforma \textbf{Github}.
    \item Aplicar la metodología ágil \textbf{Scrum} para el desarrollo del proyecto.
    \item Utilizar herramientas de integración continua \textbf{CI/CD}, como las integradas en \textbf{Github actions}, comprobando el despliege del proyecto, ejecución de test e integrando test de cobertura de código con \textbf{coverage}, control de calidad de código integrando \textbf{SonarCloud} con Github actions.
    \item Utilizar \textbf{SonarCloud} para la gestión de la calidad del código y su integración con \textbf{Github actions}
\end{itemize}

\section{Objetivos personales}\label{objetivos_personasles}
\begin{itemize}
    \item Utilizar el \textit{framework} \textbf{Django} para el desarrollo de aplicaciones web.
    \item Investigar el comportamiento de los modelos con diferentes configuraciones para conocer la viabilidad del proyecto.
    \item Conocer y utilizar \textbf{Docker} para la creación de contenedores y facilitar el despliegue del proyecto.
    \item Aprender a utilizar herramientas de control de versiones como \textbf{GIT} y plataformas como \textbf{Github}.
    \item Aprender a utilizar herramientas de integración continua como \textbf{Github actions}.
    \item Aprender a utilizar herramientas de gestión de proyectos como \textbf{Zube}.
    \item Aprender a utilizar herramientas de calidad de código como \textbf{SonarCloud}.
    \item Aprender a utilizar herramientas de documentación como \textbf{Sphinx}.
\end{itemize}
\capitulo{3}{Conceptos teóricos}


\section{Procesamiento de lenguaje natural (NLP)}

El procesamiento del lenguaje natural (NLP) es un campo de la inteligencia artificial centrado en 
la interacción entre las computadoras y el lenguaje humano que permite a los ordenadores entender, 
interpretar y trabajar con este lenguaje de forma eficaz. Esta tecnología es esencial para 
analizar información escrita o hablada, superando desafíos como las variaciones en los dialectos, 
el uso de jerga o las irregularidades gramaticales propias del lenguaje cotidiano.


A lo largo de la historia, el lenguaje natural ha evolucionado desde los primeros sistemas de 
procesamiento de texto hasta los modelos avanzados de hoy en día.

En sus inicios, el procesamiento del lenguaje natural se centraba en reglas gramaticales 
catalogando reglas sintácticas y semánticas que aumentaban en complejidad con el tiempo, 
resultando difíciles de mantener. Estos sistemas eran limitados y no podían adaptarse a la 
diversidad del lenguaje humano.

Con el tiempo, se introdujeron enfoques estadísticos basados en modelos probabilísticos que 
utilizaban grandes cantidades de datos para aprender patrones en el lenguaje. 

La siguiente etapa fue la llegada de los modelos de aprendizaje profundo \textit{deep learning}, 
que revolucionaron el campo al permitir la creación de representaciones vectoriales del lenguaje. 
Estos modelos, como las redes neuronales recurrentes (RNN) y las redes neuronales convolucionales 
(CNN), mejoraron significativamente la precisión en tareas como la traducción automática y el 
análisis de sentimientos.

El avance más reciente ha sido el desarrollo de la arquitectura Transformer. Esta tecnología 
reemplazó a las RNN tradicionales gracias a su uso del mecanismo de atención, que permite procesar
 secuencias de texto de manera más eficiente y con mayor contexto. Los Transformers emplean 
 componentes llamados \textit{encoders} y \textit{decoders} para transformar entradas en salidas, y 
 son la base de los modelos de lenguaje más avanzados utilizados actualmente, como GPT, BERT o T5.

\section{Modelos de lenguaje grandes (LLM)}




\section{Secciones}

Las secciones se incluyen con el comando section.

\subsubsection{Subsubsecciones}

Y subsecciones. 


\section{Referencias}

Las referencias se incluyen en el texto usando cite~\cite{wiki:latex}. Para citar webs, artículos o libros~\cite{koza92}, si se desean citar más de uno en el mismo lugar~\cite{bortolot2005, koza92}.


\section{Imágenes}

Se pueden incluir imágenes con los comandos standard de \LaTeX, pero esta plantilla dispone de comandos propios como por ejemplo el siguiente:

\imagen{escudoInfor}{Autómata para una expresión vacía}{.5}



\section{Listas de items}

Existen tres posibilidades:

\begin{itemize}
	\item primer item.
	\item segundo item.
\end{itemize}

\begin{enumerate}
	\item primer item.
	\item segundo item.
\end{enumerate}

\begin{description}
	\item[Primer item] más información sobre el primer item.
	\item[Segundo item] más información sobre el segundo item.
\end{description}
	
\begin{itemize}
\item 
\end{itemize}

\section{Tablas}

Igualmente se pueden usar los comandos específicos de \LaTeX o bien usar alguno de los comandos de la plantilla.

\tablaSmall{Herramientas y tecnologías utilizadas en cada parte del proyecto}{l c c c c}{herramientasportipodeuso}
{ \multicolumn{1}{l}{Herramientas} & App AngularJS & API REST & BD & Memoria \\}{ 
HTML5 & X & & &\\
CSS3 & X & & &\\
BOOTSTRAP & X & & &\\
JavaScript & X & & &\\
AngularJS & X & & &\\
Bower & X & & &\\
PHP & & X & &\\
Karma + Jasmine & X & & &\\
Slim framework & & X & &\\
Idiorm & & X & &\\
Composer & & X & &\\
JSON & X & X & &\\
PhpStorm & X & X & &\\
MySQL & & & X &\\
PhpMyAdmin & & & X &\\
Git + BitBucket & X & X & X & X\\
Mik\TeX{} & & & & X\\
\TeX{}Maker & & & & X\\
Astah & & & & X\\
Balsamiq Mockups & X & & &\\
VersionOne & X & X & X & X\\
} 

\capitulo{4}{Técnicas y herramientas}

\section{Metodologías}

\subsection{Scrum}

Scrum es un marco de trabajo ágil, líder en el desarrollo de software y la gestión de proyectos, 
que se distingue por su flexibilidad. Organiza el trabajo en sprints, iteraciones cortas de una a cuatro 
semanas, cada una entregando un incremento funcional del producto. 
Al final de cada sprint, se evalúa el resultado para detectar mejoras 
y ajustar prioridades según las necesidades del proyecto. Los requisitos, que pueden cambiar durante el 
desarrollo, se recogen en la pila del producto, mientras que las tareas específicas de cada iteración se 
detallan en la pila del sprint. Este enfoque incremental permite solapar fases del proyecto, optimizando 
tiempos y facilitando la adaptación a nuevos requerimientos. Scrum promueve la calidad a través del
aprendizaje continuo del equipo, generando entregas tempranas de productos utilizables y asegurando una 
mejora constante alineada con las necesidades del usuario~\cite{wiki:Scrum}.

\subsection{Kanban}
Kanban es un método ágil que se centra en la visualización del flujo de trabajo y la gestión del
trabajo en curso. Utiliza un tablero representando las tareas en una serie de tarjetas, permitiendo a los
miembros del equipo ver el progreso y las prioridades del trabajo.
Kanban Se basa en principios como la limitación del trabajo en curso, la mejora continua y la
adaptación a los cambios, es flexible y se puede aplicar a diferentes tipos de proyectos y equipos,
lo que lo convierte en una herramienta valiosa para la gestión de proyectos en entornos ágiles~\cite{wiki:Kanban}.

\subsection{Scrumban}
Scrumban es una metodología híbrida que combina elementos de Scrum y Kanban.
Se basa en la estructura de Scrum, con sprints y roles definidos, pero incorpora la flexibilidad de Kanban
en la gestión del flujo de trabajo.
Esta metodología es especialmente útil para equipos que buscan una transición suave entre Scrum y Kanban,
o para aquellos que desean aprovechar lo mejor de ambos enfoques.
Scrumban se centra en la mejora continua, la visualización del trabajo y la colaboración del equipo,
lo que lo convierte en una opción popular para la gestión de proyectos ágiles~\cite{web:scrumban}.

\capitulo{5}{Aspectos relevantes del desarrollo del proyecto}

\section{Inicio del proyecto}

El desarrollo de este proyecto surgió tras un proceso de negociación y diálogo con mi tutor, 
durante el cual le presenté varias propuestas iniciales. A partir de sus ideas y sugerencias, 
acordamos centrar el trabajo en la creación de un sistema para la corrección automática de 
tareas de alumnos utilizando inteligencia artificial.

La motivación principal radica en la creciente integración de tecnologías de inteligencia artificial 
y modelos de lenguaje a gran escala (\textit{LLMs}) en el ámbito educativo. Dado que el alumnado empieza a 
ayudarse de la guía proporcionada por este tipo de herramientas en la resolución de sus tareas, 
resulta natural y necesario explorar métodos que permitan 
evaluar dichas tareas utilizando la misma metodología basada en \textit{IA}.

El proyecto se seleccionó por ser una propuesta innovadora, que aprovecha tecnologías actuales 
como los \textit{LLMs} para implementar una aplicación web que se distingue de las soluciones tradicionales 
de corrección automática. Esta aplicación no solo integra inteligencia artificial generativa, sino que 
también ofrece una interfaz y funcionalidades diferentes a lo habitual, adaptándose a las necesidades 
actuales del entorno educativo digital.

Este enfoque no solo pretende agilizar el proceso de evaluación, sino también adaptarse a las nuevas 
formas de aprendizaje y producción de contenidos, manteniendo la coherencia con las herramientas que 
los alumnos emplean. De esta forma, el proyecto se posiciona como una respuesta innovadora y pertinente 
a la evolución tecnológica en el contexto académico.

\section{Gestión y metodología del proyecto}

Desde el inicio del proyecto se decidió aplicar una \textbf{metodología ágil} adaptada al contexto académico e individual del trabajo. 
En concreto, se optó por \textbf{\textit{Scrumban}}~\cite{web:scrumban}, una combinación entre \textit{Scrum} y \textit{Kanban} que permite mantener ciclos de trabajo planificados 
(\textit{sprints}) al mismo tiempo que se ofrece una mayor flexibilidad en la gestión de tareas.

El proyecto se estructuró en \textit{sprints} de dos semanas, al final de los cuales se realizaban \textbf{reuniones con el tutor}. 
En estas sesiones se revisaban las \textbf{tareas completadas} durante el \textit{sprint} anterior, 
se evaluaba el cumplimiento de los objetivos y se proponían nuevas tareas para el siguiente ciclo. 
Este enfoque iterativo facilitó una \textbf{mejora continua} del producto y permitió adaptarse a cambios o mejoras surgidas a lo largo del desarrollo.

La \textbf{gestión de tareas} se realizó con la herramienta \textbf{\emph{Zube}}, integrada con \textit{GitHub}. Esta plataforma permitió organizar el trabajo mediante 
\textit{issues}, que simulaban la dinámica de un equipo multidisciplinar. Se utilizó un \textbf{tablero \emph{Kanban}} con columnas como 
\textit{inbox}, \textit{ready}, \textit{in progress}, \textit{in review} y \textit{done}, lo que ayudó a tener una visión clara y actualizada del estado de cada tarea y a 
gestionar eficazmente el flujo de trabajo.

También se utilizaron \textbf{gráficos \textit{burndown}} para medir el progreso de cada \textit{sprint} y visualizar la carga de trabajo restante. 
En general, se logró cumplir con los plazos previstos. Estas herramientas resultaron útiles para mantener un \textbf{control claro del avance} y asegurar que los objetivos establecidos al inicio de cada ciclo se cumplieran de forma efectiva.

Adicionalmente, para el \textbf{control de versiones} se siguió la \textbf{metodología \textit{GitFlow}}, que permitió organizar el repositorio de forma estructurada 
mediante ramas diferenciadas como \textit{main}, \textit{develop}, \textit{feature}, \textit{release} y \textit{hotfix}. Esta estrategia facilitó un desarrollo más limpio, 
con mayor trazabilidad de los cambios y mejor manejo de nuevas funcionalidades y correcciones de errores.

\section{Formación}

Este trabajo ha requerido adquirir ciertos conocimientos específicos que inicialmente no se dominaban en profundidad, especialmente en lo relativo a la integración de modelos de lenguaje (\textit{LLMs}) con aplicaciones web, la ejecución de tareas en segundo plano y el despliegue de servicios mediante contenedores. No obstante, ya se contaba con una base previa en bases de datos y programación en Python, lo cual ha facilitado algunas de las fases del proyecto.

La tecnología central del \textit{backend} ha sido el \textit{framework Django}, cuyo sistema de vistas, \textit{ORM} y gestión de rutas se ha aprendido principalmente a través de su documentación oficial~\cite{web:django,web:django-orm}, y mediante pruebas experimentales con pequeños proyectos de ejemplo. La comprensión de los distintos motores de bases de datos también fue necesaria; se optó finalmente por \textit{PostgreSQL} como solución~\cite{web:postgresql}.

Uno de los mayores retos ha sido la integración de modelos de lenguaje ejecutados localmente a través de \textit{Ollama}~\cite{web:ollama}. Para ello, fue clave familiarizarse con el funcionamiento de la biblioteca \textit{LangChain}~\cite{web:langchain}, así como su extensión para \textit{Ollama}~\cite{web:langchain-ollama}.

En cuanto a la ejecución de tareas en segundo plano, fundamental para no bloquear la experiencia del usuario, se investigaron varias alternativas. Finalmente se eligió \texttt{Celery} con \texttt{Redis} como broker~\cite{web:celery,web:redis}.

Otro aspecto relevante del aprendizaje ha sido la contenerización de servicios mediante \textit{Docker} y \textit{Docker Compose}, lo cual permitió simular entornos de producción de forma local. Aunque no se tenía experiencia previa con estas tecnologías, su curva de aprendizaje fue razonable gracias a su documentación oficial y a la gran cantidad de recursos disponibles en línea~\cite{web:docker_doc}.

En cuanto al desarrollo del \textit{frontend}, se utilizaron conocimientos previos de \texttt{HTML}~\cite{web:html}, \texttt{CSS}~\cite{web:css} y \texttt{JavaScript}~\cite{web:javascript}, complementados con el uso del framework \texttt{Tailwind CSS}~\cite{web:tailwind} y su librería de componentes Flowbite~\cite{web:flowbite}, que permitieron construir una interfaz moderna y accesible.

Por último, para la documentación técnica del proyecto se exploraron y probaron herramientas como Sphinx~\cite{web:sphinx}, su tema Book Theme~\cite{web:sphinxbooktheme} y MkDocs~\cite{web:mkdocs}. También se utilizó Markdown~\cite{web:markdown} como lenguaje de marcado ligero y herramientas visuales como Draw.io~\cite{web:drawio} y Miro~\cite{web:miro} para la elaboración de diagramas y esquemas.

A lo largo del proceso se han consultado otras muchas fuentes, tanto en forma de documentación como de vídeos~\cite{web:django-curso-youtube,web:django-curso-youtube-pildoras} y artículos técnicos, pero sin duda las anteriormente mencionadas han sido las más relevantes en cuanto a la formación adquirida durante el desarrollo de este proyecto.

\section{Desarrollo de la aplicación}

El desarrollo de la aplicación se estructuró en fases iterativas, comenzando con la implementación de las funcionalidades básicas de gestión de usuarios y navegación. La primera etapa incluyó el desarrollo del sistema de autenticación (\textit{login, logout} y registro), así como la creación de una página inicial sencilla desde la cual los usuarios podían acceder al resto de funcionalidades. Esta base permitió asegurar un entorno seguro y personalizado para cada usuario.

Posteriormente, se incorporaron las funcionalidades relacionadas con la gestión de los elementos fundamentales para la evaluación: los \textit{prompts} y las rúbricas. Se implementaron las funcionalidades necesarias para crear, consultar y eliminar prompts personalizados, así como importar rúbricas de evaluación en formato markdown, con el objetivo de facilitar la configuración de criterios de corrección.

En una siguiente iteración, se desarrolló la lógica para configurar las correcciones propiamente dichas. En esta fase, la aplicación permitió asociar un \textit{prompt}, una rúbrica y subir un fichero comprimido con las tareas de programación ubicadas en la raíz del fichero. Estas configuraciones quedaban almacenadas en la base de datos, permitiendo mantener un registro estructurado de cada tarea de evaluación preparada por el usuario.

Con la configuración establecida, se implementó el procesamiento en segundo plano de las correcciones utilizando \texttt{Celery} junto con \texttt{Redis} como \textit{message broker}. Esta solución permitió ejecutar de forma asíncrona las tareas relacionadas con el envío de información al modelo de lenguaje y la recepción de respuestas, evitando bloquear el servidor principal.

Debido a que los \textit{LLMs} pueden tardar varios segundos en generar una respuesta, especialmente en entornos locales, resultaba inviable gestionar estas operaciones de forma síncrona. Gracias a esta arquitectura, la aplicación mantiene su capacidad de respuesta y permite gestionar múltiples evaluaciones simultáneamente sin afectar al rendimiento general.

Finalmente, se añadió la posibilidad de modificar ciertos parámetros del modelo \textit{LLM} en cada evaluación, como el nivel de creatividad (\textit{temperatura}), la probabilidad acumulada máxima (\textit{top-p}) y el número máximo de \textit{tokens} considerados en la selección (\textit{top-k}). Esta flexibilidad permitió a los usuarios adaptar la evaluación a distintos escenarios y necesidades específicas.

\subsection*{Procesamiento de archivos y estructura del sistema}

Se estableció un sistema de subida de archivos donde el usuario podía cargar:
\begin{itemize}
    \item Un archivo con la rúbrica de evaluación (en formato \texttt{md}).
    \item Un archivo comprimido con las tareas en su interior.
\end{itemize}

Además, el usuario podía escribir manualmente un \textit{prompt} personalizado. La lógica del sistema se encargaba de ensamblar todos estos elementos en una única entrada textual coherente, que era enviada al modelo.

\subsection*{Evaluación con modelos LLM}

El modelo principal utilizado fue \texttt{LLaMA 3.1}, ejecutado localmente a través de \textit{Ollama} e integrado en la lógica de la aplicación mediante la librería \texttt{langchain-ollama}.

Cada evaluación se procesaba como una tarea independiente, registrando:
\begin{itemize}
    \item Fecha y hora de la configuración.
    \item \textit{Prompt} y rúbrica utilizados.
    \item Tareas evaluadas.
    \item Fecha y hora de la evaluación.
    \item Parámetros del modelo y formato de la respuesta.
    \item Respuesta devuelta por el modelo.
\end{itemize}

Cada evaluación generaba un archivo de texto plano (\texttt{.txt}) que se almacenaba en el directorio \texttt{media}.

\subsection*{Interfaz y usabilidad}

La interfaz web fue desarrollada con \texttt{HTML}, \texttt{CSS} y \texttt{JavaScript}, apoyada por \texttt{Tailwind CSS} y la biblioteca de componentes \texttt{Flowbite}, lo que permitió un diseño limpio y funcional. Se priorizó una navegación simple en la que el flujo de trabajo fuera evidente: autenticación → gestión de prompts  → gestión de rúbricas → configuración de evaluación → ejecución de correcciones → consulta o descarga de resultado.

Se incorporaron funcionalidades adicionales como:
\begin{itemize}
    \item Tabla de correcciones por usuario.
    \item Tabla de rúbricas por usuario.
    \item Tabla de prompts por usuario.
    \item Gestión de usuarios y sesiones mediante el sistema de autenticación de Django.
\end{itemize}

\subsection*{Contenerización y despliegue}

Para garantizar la reproducibilidad y facilitar el despliegue en distintos entornos, se optó por contenerizar todos los servicios mediante \texttt{Docker}. El sistema quedó estructurado en varios contenedores:
\begin{itemize}
    \item \textit{Django}.
    \item Base de datos \textit{PostgreSQL}.
    \item \textit{Broker} de mensajería \textit{Redis}.
    \item \textit{Ollama} con modelos \textit{LLM}.
    \item \textit{Celery}.
\end{itemize}

Esta separación de servicios permitió trabajar de forma más ordenada y facilitó el proceso de pruebas y desarrollo local.

\subsection*{Resumen de funcionalidades desarrolladas}

A lo largo del desarrollo se implementaron las siguientes características clave:

\begin{itemize}
    \item Sistema de autenticación con \textit{login, logout} y registro.
    \item Gestión y edición de \textit{prompts} y rúbricas.
    \item Configuración de tareas de evaluación que almacenan \textit{prompt}, rúbrica y archivos asociados.
    \item Ejecución asincrónica de evaluaciones mediante Celery y Redis.
    \item Personalización de parámetros del modelo \textit{LLM} en cada evaluación.
    \item Registro, almacenamiento y consulta de resultados.
    \item Interfaz web amigable y funcional.
    \item Contenerización completa para facilitar el despliegue.
\end{itemize}

El desarrollo modular y progresivo permitió construir una aplicación fácilmente ampliable, alineada con los objetivos establecidos al inicio del proyecto.

\subsection{Constantes configurables del sistema}

El sistema cuenta con ciertas constantes definidas en el código fuente que permiten ajustar su comportamiento de forma sencilla. En particular, existen dos elementos clave que pueden modificarse para ampliar la funcionalidad de la aplicación:

\begin{itemize}
    \item En el archivo \texttt{corrections/views.py}, línea 28, se encuentra la tupla \texttt{VALID\_EXTENSION}, que especifica las extensiones de archivo válidas para la corrección automática. Actualmente incluye formatos como \texttt{.java}, \texttt{.sql}, \texttt{.py}, \texttt{.cpp}, \texttt{.txt}, entre otros. Esta lista puede ampliarse fácilmente para soportar nuevos lenguajes o tipos de archivo.
    
    \item En el archivo \texttt{corrections/forms.py}, línea 10, se define la lista \texttt{MODEL\_CONTEXT\_CHOICES}, que establece las opciones de contexto máximo en número de \textit{tokens} que pueden procesarse por el modelo de lenguaje. Los valores actuales son 2048, 4096 y 8192 \textit{tokens}, pero esta lista puede modificarse para incorporar futuras configuraciones.
\end{itemize}

Estas constantes permiten adaptar la aplicación a distintos entornos o requerimientos sin necesidad de modificar la lógica principal del sistema.


\section{Problemas, incidencias y soluciones}

Durante el desarrollo de la aplicación surgieron distintos problemas técnicos, principalmente relacionados con la ejecución de tareas asíncronas mediante \texttt{Celery} y la contenerización del proyecto con \texttt{Docker}. A continuación, se describen los problemas más relevantes y las soluciones aplicadas.

\subsection*{Inicio con \texttt{async} de \textit{Django} y cambio a \texttt{Celery} y \texttt{Redis}}

En una primera etapa, se intentó utilizar las funcionalidades \texttt{async} nativas de \textit{Django} para gestionar la ejecución de tareas en segundo plano. Sin embargo, esta implementación inicial no cumplía con las necesidades reales de procesamiento asíncrono, ya que no permitía la ejecución fuera del ciclo de vida de la solicitud \textit{HTTP} ni garantizaba la persistencia ni la gestión eficiente de tareas largas.

\textbf{Solución:} Se optó por integrar \texttt{Celery} como sistema de cola de tareas junto con \texttt{Redis} como broker de mensajes. Esta combinación permitió una gestión robusta y confiable de las tareas asíncronas, garantizando que las operaciones de evaluación y generación de respuestas pudieran ejecutarse en segundo plano sin afectar la experiencia del usuario.

\subsection*{Advertencia por ejecutar Celery como usuario \texttt{root}}

Uno de los primeros problemas detectados fue una advertencia al iniciar \texttt{Celery}, que alertaba de que el proceso se estaba ejecutando como usuario \texttt{root}. Aunque funcional, esto supone un riesgo potencial de seguridad, especialmente en entornos de producción, y va en contra de las buenas prácticas recomendadas.

\textbf{Solución:} Para evitar ejecutar procesos como superusuario, se creó un nuevo usuario dentro del contenedor de Docker (llamado \texttt{dj\_admin}) y se configuró para que los procesos de Django y Celery se ejecutaran bajo ese usuario. En el archivo \texttt{Dockerfile}, se utilizó la instrucción \texttt{USER dj\_admin} para establecer este comportamiento. Esto permitió mantener un entorno más seguro y alineado con las recomendaciones de los \textit{frameworks} utilizados.

\subsection*{Problemas de permisos al escribir en la carpeta \texttt{media}}

Tras aplicar la solución anterior, surgió una nueva incidencia: el usuario \texttt{dj\_admin} no tenía permisos suficientes para acceder a la carpeta \texttt{media}, que se utiliza para almacenar archivos temporales y resultados de las evaluaciones (en formato \texttt{.txt}). Esta carpeta es esencial para el correcto funcionamiento del sistema, ya que en ella se guardan las tareas a corregir, así como la respuesta generada por el modelo \textit{LLM}.

\textbf{Solución:} Se añadieron comandos específicos al \texttt{Dockerfile} para asignar los permisos adecuados a la carpeta \texttt{/app/media}. Concretamente, se utilizó el siguiente bloque de instrucciones:

\begin{verbatim}
...

RUN chown -R dj_admin:dj_admin /app/media

EXPOSE 8000

COPY entrypoint.sh /entrypoint.sh

RUN chmod +x /entrypoint.sh

ENTRYPOINT ["/entrypoint.sh"]

USER dj_admin

CMD ["python", "manage.py", "runserver", "0.0.0.0:8000"]

\end{verbatim}

Esto garantiza que el usuario sin privilegios tenga acceso completo a dicha carpeta. Con este ajuste, fue posible evitar errores de acceso y escritura al ejecutar tareas de evaluación que involucran la generación y almacenamiento de archivos.

\subsection*{Conflictos por carpetas \texttt{media} durante la ejecución de tests}

Durante la ejecución de los \textit{tests} automatizados de la aplicación de correcciones, se detectó que se creaban carpetas nuevas dentro del directorio \texttt{media}. Esto podía provocar interferencias o conflictos con datos y archivos pertenecientes a otros usuarios, afectando la integridad y el aislamiento de las pruebas.

\textbf{Solución:} Para evitar estos problemas, se implementó un mecanismo que crea una carpeta temporal específica para \texttt{media} cada vez que se ejecutan los \textit{tests}. Esta carpeta temporal se utiliza exclusivamente durante la ejecución de las pruebas y se elimina automáticamente al finalizar, garantizando así que el entorno de \textit{tests} esté aislado y no interfiera con el sistema en uso normal.

Esta estrategia asegura un entorno limpio y reproducible para las pruebas, mejorando la fiabilidad y evitando efectos colaterales no deseados en el almacenamiento de archivos durante la fase de \textit{testing}.

\capitulo{6}{Trabajos relacionados}

\subsection{Herramienta de corrección automática de prácticas de programación}

El trabajo \textit{Herramienta de corrección automática de prácticas de programación}~\cite{tfg:alvarez2024}, realizado en la Universidad Politécnica de Madrid, presenta un sistema para la corrección automática de ejercicios de programación. La herramienta analiza el código de los estudiantes, teniendo en cuenta aspectos como el uso de buenas prácticas, la legibilidad y la estructura del código, con el objetivo de ofrecer retroalimentación inmediata y aliviar la carga de trabajo del profesorado.

A diferencia de IAGScore, que utiliza modelos de lenguaje a gran escala (\textit{LLMs}) para interpretar y evaluar el contenido de las tareas, este sistema se basa en reglas definidas previamente y en análisis más estáticos. Aunque ambas propuestas buscan automatizar la corrección en cursos de programación, lo hacen con metodologías distintas: una más tradicional, basada en validaciones estructurales, y otra más flexible, apoyada en inteligencia artificial generativa.

\subsection{Sistema de corrección automática basado en tests y análisis estático}

El trabajo presentado como \textit{Sistema de corrección automática basado en tests y análisis estático}~\cite{montanes2014} describe un sistema para la corrección automática de tareas de programación que se apoya principalmente en la ejecución de tests unitarios para verificar que el código entregado cumple con los requisitos funcionales. Además, incorpora técnicas de análisis estático para evaluar aspectos como el estilo, la calidad y la estructura del código.

Este enfoque combina la validación dinámica mediante pruebas automatizadas con la revisión estática del código, ofreciendo una evaluación más completa que se centra tanto en el funcionamiento correcto como en las buenas prácticas de programación. 

A diferencia del sistema propuesto en este proyecto, que emplea modelos de lenguaje (\textit{LLMs}) para interpretar y evaluar las soluciones desde una perspectiva semántica y contextual, el enfoque descrito se basa en validaciones formales y reglas predefinidas. Mientras que la herramienta basada en tests requiere una programación previa de casos de prueba y métricas específicas, el uso de \textit{LLMs} permite una evaluación más flexible y adaptable al tipo de problema planteado, así como una retroalimentación en lenguaje natural más cercana a la corrección humana.

\subsection{Rubric Is All You Need}

El artículo \textit{Rubric Is All You Need: Enhancing LLM-based Code Evaluation With Question-Specific Rubrics}~\cite{pathak2025}, analiza el uso de modelos de lenguaje a gran escala para evaluar tareas de programación en cursos universitarios. Proponen tres técnicas específicas: Evaluación Completa de Rúbrica (CRE), Evaluación Punto a Punto (PRE) y Evaluación por Ensamble (EME), que permiten evaluar la lógica y el contenido semántico del código más allá de la simple corrección sintáctica. 

El estudio se basa en dos conjuntos de datos: exámenes de Java en Programación Orientada a Objetos y ejercicios de estructuras de datos y algoritmos. Los resultados muestran que el uso de rúbricas específicas para cada problema mejora la precisión y la calidad de la retroalimentación, obteniendo una alta correlación con las calificaciones de evaluadores humanos. También introducen una métrica llamada \textit{Leniency}, que ajusta el nivel de exigencia del sistema evaluador, aportando flexibilidad para adaptar la rigurosidad de la corrección, un aspecto interesante para futuras mejoras del proyecto, especialmente en entornos educativos con distintos niveles de exigencia.

\small
\tabla{Comparativa entre trabajos de corrección automática de código e IAGScore}
{p{2.8cm} p{2cm} p{2cm} p{2cm} p{2cm}}{5}{comparativaTrabajosBinaria}
{
\textbf{Características} & 
\textbf{Corrección automática (UPM)} & 
\textbf{Tests y análisis estático} & 
\textbf{Rubric Is All You Need} & 
\textbf{IAGScore} \\
}
{
Evaluación automatizada & Sí & Sí & Sí & Sí \\
Análisis estático & Sí & Sí & No & No \\
Uso de rúbricas & No & Parcial & Sí & Sí \\
Generación de lenguaje natural & No & No & Sí & Sí \\
Uso de IA/LLM & No & No & Sí & Sí \\
Evaluación contextual & No & No & Parcial & Sí \\
Adaptabilidad & Media & Baja & Alta & Alta \\
Personalización & No & No & Parcial & Sí \\
}
\normalsize

\capitulo{7}{Conclusiones y Líneas de trabajo futuras}

\section{Conclusiones}

A lo largo del desarrollo de este Trabajo de Fin de Grado se han alcanzado los objetivos propuestos, tanto a nivel funcional como técnico y personal. El proyecto ha consistido en la construcción de una herramienta web capaz de evaluar automáticamente ejercicios de programación, integrando modelos de lenguaje a gran escala (\textit{LLMs}) de forma local.

\subsection{Cumplimiento de los objetivos}

Se ha logrado integrar modelos \textit{LLM} de forma local a través de \textit{Ollama}. Esto ha permitido implementar un sistema de corrección automatizada capaz de analizar el código enviado por el usuario y devolver una evaluación textual. Además, se ha desarrollado una interfaz que permite interactuar con el sistema de manera sencilla.

Durante el desarrollo se han utilizado tecnologías como \textit{Django, PostgreSQL, Docker, GitHub Actions o Sphinx}, entre otras. El uso de contenedores ha facilitado el despliegue y prueba del sistema, y el control de versiones ha permitido mantener una estructura organizada. También se han incorporado prácticas de calidad del software, como la integración de \textit{SonarCloud}.

Este proyecto ha sido una oportunidad para afianzar conocimientos en múltiples áreas. Se han reforzado competencias en gestión de proyectos, control de versiones, calidad de código y documentación técnica, permitiendo una visión más completa del ciclo de vida del \textit{software}. 

\subsection{Reflexiones y retos del proceso}

Uno de los principales desafíos ha sido trabajar con tecnologías emergentes como los modelos de lenguaje de gran escala (\textit{LLMs}). Desde el inicio del proyecto, existía incertidumbre sobre su viabilidad para cumplir los objetivos planteados, lo que supuso un reto tanto en el diseño como en la implementación. Esta falta de garantías iniciales exigió una actitud exploratoria, así como la adaptación continua a un entorno tecnológico en rápida evolución.

Además, se ha tenido que lidiar con algunas limitaciones inherentes a este tipo de modelos, como las denominadas \textit{alucinaciones}, es decir, respuestas generadas que pueden parecer coherentes pero no son correctas o están inventadas. Dado que se trata de tecnologías emergentes, este fenómeno subraya la necesidad de mantener una supervisión humana en el proceso de evaluación, especialmente en entornos educativos donde la precisión y la fiabilidad son fundamentales.

\subsection{Valoración personal}

La realización de este proyecto ha sido altamente satisfactoria, tanto a nivel académico como personal. No solo se ha desarrollado una aplicación funcional, sino que se ha conseguido una aproximación a la integración de \textit{LLMs} en el entorno educativo para la corrección de tareas. 

Esta experiencia ha sido fundamental tanto para consolidar la formación académica como para prepararse para futuros retos profesionales.

\section{Líneas de trabajo futuras}

El proyecto desarrollado sienta las bases para futuras ampliaciones tanto a nivel funcional como de integración con otros entornos. Una de las líneas de trabajo más prometedoras es la implementación de una \textit{API} que permita conectar el sistema con plataformas educativas como \textit{Moodle}. Esta integración facilitaría su adopción en entornos reales de enseñanza, automatizando la evaluación de tareas de programación directamente desde el entorno de gestión académica.

Otra posible ampliación consiste en utilizar distintos modelos \textit{LLM} para realizar una evaluación cruzada del resultado generado por el \textit{LLM} evaluador principal. Esto podría mejorar la precisión y robustez del sistema, al aprovechar distintas perspectivas y reducir posibles sesgos del modelo individual.

Asimismo, resultaría adecuado implementar un sistema de notificaciones que informe al usuario una vez finalizado el proceso de corrección. Estas notificaciones podrían realizarse mediante correo electrónico, \textit{SMS} o notificaciones \textit{push}, mejorando así la experiencia de usuario y facilitando la gestión del tiempo de espera en procesos asíncronos.


\bibliographystyle{plain}
\bibliography{bibliografia}

\end{document}
